\chapter{Charakterisierungen und Randfälle}
\label{b04guide:ch:criteria}

\section{Wie läßt sich Totalität erkennen?}

\begin{theorem}[Faserkriterium]\label{b04guide:thm:fiber}
Für $R\subseteq A\times B$ sind äquivalent:
\[
 \begin{aligned}
 &\TR(R;A,B),\\
 &\forall a\in A\,(R_a\ne\varnothing),\\
 &\{a\in A\mid\exists b\in B\,((a,b)\in R)\}=A.
 \end{aligned}
\]
Die letzte Gleichheit besagt, daß jedes Element der
vorgegebenen Ausgangsmenge tatsächlich vorkommt.
\end{theorem}

\begin{theorem}[Nichtleere Teilmengen erreichen etwas]
\label{b04guide:thm:nonempty}
Für $R\subseteq A\times B$ sind äquivalent:
\[
 \begin{aligned}
 &\TR(R;A,B),\\
 &\forall C\subseteq A\,
          (C\ne\varnothing\Rightarrow R[C]\ne\varnothing),\\
 &\forall C\subseteq A\,
          (R[C]=\varnothing\Longleftrightarrow C=\varnothing).
 \end{aligned}
\]
\end{theorem}

\begin{theorem}[Lokales Kriterium auf einer Überdeckung]
\label{b04guide:thm:local}
Seien $\mathcal C\subseteq\mathcal P(A)$,
$\bigcup\mathcal C=A$ und $R\subseteq A\times B$. Dann gilt
\[
 \TR(R;A,B)\ \Longleftrightarrow\
 \forall C\in\mathcal C\,
       \TR\bigl(R\cap(C\times B);C,B\bigr).
\]
Die überdeckenden Teilmengen müssen nicht disjunkt sein.
Die leere Überdeckung ist für $A=\varnothing$ zugelassen.
\end{theorem}

\begin{theorem}[Ausgangsmenge und Zielmenge]
\label{b04guide:thm:carriers}
Ist $\TR(R;A,B)$, so ist die Ausgangsmenge durch $R$ bestimmt:
\[
 A=\{a\in\bigcup\bigcup R\mid
                  \exists b\,((a,b)\in R)\}.
\]
Hier sind die geordneten Paare wie in Band~3 als
$(a,b)=\{\{a\},\{a,b\}\}$ dargestellt.
Für jede Menge $D$ gilt außerdem
\[
 \TR(R;A,D)\quad\Longleftrightarrow\quad R[A]\subseteq D.
\]
Insbesondere ist $R[A]$ die kleinste mögliche Zielmenge
derselben totalen Relation; die Zielmenge ist im allgemeinen
nicht eindeutig.
\end{theorem}

\section{Leere Mengen und Existenz}

\begin{theorem}[Leere Ausgangsmenge und leere Zielmenge]
\label{b04guide:thm:empty}
Für alle Mengen $A,B$ gilt
\[
 \begin{aligned}
 \TRset(\varnothing,B)&=\{\varnothing\},\\
 \TRset(A,\varnothing)&=
 \begin{cases}
   \{\varnothing\},& A=\varnothing,\\
   \varnothing,& A\ne\varnothing.
 \end{cases}
 \end{aligned}
\]
Für $\TR(R;A,B)$ gilt ferner
\[
 R=\varnothing\Longleftrightarrow A=\varnothing.
\]
Ist $A\ne\varnothing$, so sind auch $R[A]$ und $B$ nichtleer.
\end{theorem}

\begin{theorem}[Existenz und volle Rechtecke]
\label{b04guide:thm:existence}
Für alle Mengen $A,B$ sind äquivalent:
\[
 \begin{aligned}
 &\TRset(A,B)\ne\varnothing,\\
 &A=\varnothing\ \lor\ B\ne\varnothing,\\
 &\TR(A\times B;A,B).
 \end{aligned}
\]
Wenn totale Relationen von $A$ nach $B$ existieren, enthält
$A\times B$ jede von ihnen und ist selbst total.
\end{theorem}

\begin{theorem}[Ein vorgegebenes gemeinsames Zielelement]
\label{b04guide:thm:constant}
Seien $b\in B$ und $R\subseteq A\times B$. Dann gilt
\[
 A\times\{b\}\subseteq R\quad\Rightarrow\quad \TR(R;A,B).
\]
Jede Menge $A$ besitzt also eine totale Relation nach jeder
nichtleeren Menge $B$, deren sämtliche Fasern ein gemeinsames
Element enthalten.
\end{theorem}
